\documentclass[../DoAn.tex]{subfiles}
\begin{document}
\section{Kết luận}

Trong khuôn khổ đồ án tốt nghiệp 1, sinh viên đã xây dựng thành công một hệ thống chatbot hỗ trợ giảng viên tự động trả lời các câu hỏi đơn giản của sinh viên thông qua việc fine-tune mô hình ngôn ngữ TinyLLaMA bằng kỹ thuật LoRA trên dữ liệu email thực tế. Đồ án đã giải quyết được các vấn đề chính sau:

\begin{itemize}
    \item Xây dựng pipeline hoàn chỉnh từ thu thập email, tiền xử lý, phân loại thành cặp hỏi – đáp, đến huấn luyện mô hình và sinh phản hồi.
    \item Ứng dụng mô hình nhẹ TinyLLaMA kết hợp với phương pháp fine-tune hiệu quả LoRA để huấn luyện trên tập dữ liệu nhỏ trong điều kiện tài nguyên hạn chế.
    \item Triển khai đánh giá mô hình thông qua thực nghiệm định tính và định lượng, cho thấy chatbot có khả năng phản hồi hợp lý và phù hợp với ngữ cảnh học vụ.
\end{itemize}

Tuy nhiên, đồ án vẫn còn tồn tại một số hạn chế:

\begin{itemize}
    \item Mô hình chưa xử lý tốt các câu hỏi ngắn, thiếu ngữ cảnh hoặc dạng mới lạ chưa từng xuất hiện trong dữ liệu huấn luyện.
    \item Chưa tích hợp khả năng truy xuất dữ liệu thời gian thực như lịch học, thông báo mới từ hệ thống học vụ.
    \item Đánh giá mô hình mới dừng ở mức độ cơ bản, chưa sử dụng tập benchmark chuẩn hay áp dụng phương pháp đánh giá ngữ nghĩa nâng cao.
\end{itemize}

\section{Hướng phát triển trong tương lai}

Dựa trên kết quả đạt được, một số hướng phát triển khả thi trong tương lai bao gồm:

\begin{itemize}
    \item Mở rộng tập dữ liệu huấn luyện với các nguồn khác như câu hỏi từ diễn đàn học tập, LMS, chatbot sẵn có để tăng tính đa dạng.
    \item Áp dụng các kỹ thuật RAG (Retrieval-Augmented Generation) kết hợp Knowledge Graph để tăng độ chính xác, tính mạch lạc và khả năng suy luận của chatbot.
    \item Tích hợp giao diện trực quan để giảng viên và sinh viên tương tác trực tiếp với chatbot, đồng thời thu thập phản hồi để cải tiến chất lượng mô hình.
    \item Nghiên cứu sử dụng các mô hình ngôn ngữ lớn hơn (LLaMA 3, Mistral, Phi-2,...) kết hợp với LoRA hoặc QLoRA để cải thiện chất lượng mà vẫn tiết kiệm chi phí tính toán.
\end{itemize}

Tổng kết lại, đồ án là bước khởi đầu tốt cho việc ứng dụng NLP và mô hình ngôn ngữ nhỏ trong môi trường giáo dục, mở ra nhiều tiềm năng ứng dụng thực tiễn trong hỗ trợ học vụ tự động.

\end{document}